Notes:

In my solution to this problem, I actually use an advanced theorem from number theory. The theorem says that $$\pi(n) < \frac{1.25506n}{\ln(n)},$$ where $\pi(n)$ is the prime-counting function, which denotes the number of primes less than or equal to $n$. The theorem is supported by proof; it was proven in 1962 by Rosser and Schoenfeld.

You can find their proof at the following websites:

\begin{enumerate}
\item https://projecteuclid.org/journals/illinois-journal-of-mathematics/volume-6/issue-1/Approximate-formulas-for-some-functions-of-prime-numbers/10.1215/ijm/1255631807.full
\item http://denise.vella.chemla.free.fr/Rosser-Schoenfeld-1962.pdf
\item https://scispace.com/pdf/approximate-formulas-for-some-functions-of-prime-numbers-c2u470dp6o.pdf
\end{enumerate}

I thought it would help to cite them; otherwise, it looks like I'm waving a magic wand. Number theory is a rich and fascinating subject. The prime-counting function $\pi(n)$ is one of my favorite functions from number theory.

The theorem I refer to is actually corollary 1 in section 3 of the publication, on page 69.

I'm under the impression that the Rosser-Schoenfeld publication does not prove the Prime Number Theorem; instead, it uses the Prime Number Theorem as a foundation. I was able to find a proof of the Prime Number Thorem here: https://people.mpim-bonn.mpg.de/zagier/files/doi/10.2307/2975232/fulltext.pdf. I think the earliest proofs date back to the 19th century, with the work of Hadamard and de la Vallée Poussin. Gauss was able to form the conjecture in the 1790s, and he may have proven it, or come up with a lot of the ideas, himself.
